% registry: DANGVAN
\methodhead{DANGVAN}{Dang Van (1973): mesoscopic shear and hydrostatic stress}\label{sheet:DANGVAN}

\block{Source} K.~Dang Van, \emph{Sciences et Technique de l'Armement} 47 (1973) 647--722
\cite{dangvan1973}. Implemented as the documented variant of Eq.~(3) of B\"ohme, Papuga and Lange,
\emph{Fatigue Fract.\ Eng.\ Mater.\ Struct.}\ (2026) \cite{bohme2026}, the modified form of Papuga and
Halama; Eqs.~(3)--(4), (16); Table~4.

\block{Idea and scope} Mesoscopic fatigue-limit criterion. Cracks nucleate in the least favourably
oriented grains; after elastic shakedown, the mesoscopic shear stress in those grains, combined with
the hydrostatic stress, must stay below a limit line. Proposed for the high-cycle fatigue limit of
metallic components, widely used in industrial design.

\block{Inputs}
\begin{center}\small
\begin{tabular}{@{}llp{100mm}@{}}
\toprule
Symbol & Unit & Meaning\\\midrule
$s_{-1}=\fm$, $\tm$ & MPa & fully reversed tension (bending) and torsion fatigue limits\\
$\kappa$ & -- & limit ratio $\fm/\tm$ (Eq.~4)\\
$\tau(t)$ & MPa & mesoscopic shear stress (Tresca measure of the deviatoric stress minus its shakedown centre)\\
$\sigma_H(t)$ & MPa & hydrostatic stress, $\tr\Tsig/3$\\
$a$, $b$ & --, MPa & constants of the classical form, Eq.~\eqref{eq:dv-classic}\\
$\n$ & -- & unit normal of a material plane\\
$\tau_{sa}(\n)$ & MPa & radius of the smallest circle enclosing the shear-stress path on the plane $\n$\\
$\sigma_{H,a}$, $\sigma_{H,m}$, $\sigma_{H,\max}$ & MPa & amplitude, mean and maximum $\sigma_{H,m}+\sigma_{H,a}$ of $\sigma_H$\\
$\sigma_{eq,DV}$ & MPa & equivalent stress, Eq.~\eqref{eq:dv-main}\\
$\sigma_a$ & MPa & amplitude of the worked example's loads\\
\bottomrule
\end{tabular}
\end{center}

\block{Formulation} Dang Van's criterion (in the form popularised in later papers) states that, at every
instant of the stabilised cycle,
\begin{equation}
\max_t\bigl[\tau(t) + a\,\sigma_H(t)\bigr] \le b , \qquad
a = 3\Bigl(\frac{\tm}{\fm} - \frac12\Bigr) , \qquad b = \tm .
\label{eq:dv-classic}
\end{equation}
The original papers define $\tau$ through the elastic-shakedown state and give no closed form for
general loading. The repository implements Eq.~(3) of B\"ohme et al.\ (2026),
\begin{equation}
\sigma_{eq,DV} = \kappa\,\max_{\n}\tau_{sa}(\n) + \Bigl(3 - \tfrac32\kappa\Bigr)\sigma_{H,\max} ,
\qquad \kappa = \frac{\fm}{\tm} , \qquad \mathrm{DP} = \frac{\sigma_{eq,DV}}{\fm} .
\label{eq:dv-main}
\end{equation}
Dividing by $\kappa$ gives $\tau_{sa}+3(\tm/\fm-1/2)\,\sigma_{H,\max}\le\tm$, the classical constants of
Eq.~\eqref{eq:dv-classic}. For equibiaxial loading of amplitude $\sigma_a$, Eq.~(16) gives
$\sigma_{eq,DV}=2\sigma_a(1-\kappa/4)$.

\block{Constants and calibration} Paper and repository: $a$ and $b$ (or $\kappa$) from the fully reversed
tension and torsion limits, so the index is~1 at both limits. The repository uses the workbook
\texttt{Sigm1}, \texttt{Taum1}; nothing is fitted except the master-curve constants on the uniaxial
$R=-1$ design tests.

\block{Life} No native life law: Dang Van is a fatigue-limit criterion. The index is the damage parameter
of the shared master curve (Section~\ref{sec:mastercurve}).

\block{Repository choices}
\begin{itemize}
\item Amplitude plus maximum, not the maximum of the sum: the variant adds the largest shear amplitude
  and the largest hydrostatic stress instead of maximising $\tau(t)+a\,\sigma_H(t)$ over time. The two
  coincide for proportional loading. For a single harmonic the shakedown centre is the mean deviatoric
  stress, so $|\tau(t)|\le\tau_{sa}$ and, with $a>0$ ($\tm/\fm>1/2$), the variant is conservative.
\item Shear amplitude: minimum circumscribed circle on each plane (exact for a single harmonic, as in
  BP and CAIM), maximised over a sphere grid with pattern-search refinement.
\item $\sigma_{H,\max}=\sigma_{H,m}+\sigma_{H,a}$ from the harmonic.
\item Case-level model (it uses the whole history): no \texttt{\_ext} variant.
\item The index is the damage parameter of the master curve.
\end{itemize}

\block{Worked example} Synthetic material $\fm=300$, $\tm=200\MPa$ ($\kappa=1.5$). Both limits give
index~1. Equibiaxial amplitude $150\MPa$: index $0.6250$, equal to $2\sigma_a(1-\kappa/4)/\fm$. Tension
$\sigma_a=200\MPa$ with mean $100\MPa$ ($\tau_{sa}=\sigma_a/2$, $\sigma_{H,\max}=(100+200)/3\MPa$):
index $0.7500$.

\block{Verification} Eq.~(16) closed forms for equibiaxial and equitriaxial loading over several
$\kappa$; the Table~4 DV statistics of B\"ohme et al.\ (2026) on its 18 tests (maximum 20.0,
minimum $-18.6$, mean 3.6, standard deviation 11.9 points; 44.4\,\% within $\pm10$); invariance to
rotation and to a time shift. The DV column ($I^b$) of Vu, Halm and Nadot \cite{vu2010} is reproduced on
115 of 117 tests. Not implemented: the original shakedown-based formulation for arbitrary paths.
